\documentclass[10pt,a4paper]{article}
\usepackage[T1]{fontenc}
\usepackage[utf8]{inputenc}
\usepackage[brazilian]{babel}
\usepackage{lmodern}
\usepackage[margin=1.8cm]{geometry}
\usepackage{microtype}
\usepackage{amsmath,amssymb}
\usepackage{booktabs,tabularx,array,colortbl}
\usepackage{enumitem}
\usepackage{needspace}
\usepackage{pgfplots}
\usepackage{tabuleiro}
\usepackage[hidelinks]{hyperref}

\setlength{\parindent}{0pt}
\setlength{\parskip}{4pt}
\setlist{itemsep=1pt,topsep=2pt,leftmargin=1.2em}
\renewcommand{\arraystretch}{1.25}
\newcolumntype{L}{>{\raggedright\arraybackslash}p{0.24\linewidth}}
\newcolumntype{Y}{>{\raggedright\arraybackslash}X}
\definecolor{cinzaCampo}{HTML}{F1EFEA}
\definecolor{cinzaDica}{HTML}{6B7078}
\newcommand{\fn}[1]{\textup{\textsc{#1}}}
\newcommand{\campo}[2]{\cellcolor{cinzaCampo}\textbf{#1}\newline{\footnotesize\color{cinzaDica}\itshape #2}}
\newcommand{\secao}[1]{\par\needspace{14\baselineskip}\medskip{\large\bfseries\color{corTabEsc}#1}\par\vspace{2pt}}
\newcommand{\aviso}[2]{\par\smallskip\fcolorbox{corDestaque}{corDestaque!6}{\parbox{\dimexpr\linewidth-2\fboxsep-2\fboxrule}{\small\textbf{#1}\enspace #2}}\par}
\newcommand{\ok}{\textcolor{corV}{\boldmath$\checkmark$}}

\begin{document}

{\small PUCPR \enspace|\enspace Escola Politécnica \enspace|\enspace Bacharelado em Ciência da Computação\\
Inteligência Artificial (PCO118A) \enspace·\enspace Prof. Carlos Costa Cedro}

\medskip
{\LARGE\bfseries TDE2 --- Busca Competitiva}\\[2pt]
{\large Ficha de projeto \enspace·\enspace preencher em sala em 14/10}\\[2pt]
{\small RA2 \enspace·\enspace peso 4,0 \enspace·\enspace execução de 20/10 a 27/10}

\smallskip
{\small\itshape O trabalho é construir um agente para um jogo de soma zero com informação perfeita e avaliar seu
desempenho em comparação com outros agentes, que é literalmente o texto do indicador ID2.3.}

\begin{tabularx}{\linewidth}{@{}LY@{}}
\toprule
\campo{Equipe}{nomes completos} &
Guilherme Arcanjo de Morais\newline Lucas da Costa Paula\newline Renan da Silva Andrade\newline Renan Americo Herculano \\
\bottomrule
\end{tabularx}

% ---------------------------------------------------------------------
\secao{1. O jogo}
\begin{tabularx}{\linewidth}{@{}LY@{}}
\toprule
\campo{Jogo escolhido}{nome e regras em duas linhas} &
Lig 4 (Connect Four), tabuleiro vertical de 6 linhas × 7 colunas.\newline
X e O alternam soltando uma peça em uma coluna não cheia; a peça cai até a casa livre mais baixa. Vence quem alinhar
4 peças na horizontal, vertical ou diagonal; tabuleiro cheio sem alinhamento é empate. \\
\midrule
\campo{Verificação}{é de dois jogadores? de soma zero? de informação perfeita? determinístico? finito? responda as cinco} &
\textbf{Dois jogadores:} sim, X e O se alternam.\newline
\textbf{Soma zero:} sim. A vitória de um é a derrota do outro e o empate vale 0 para os dois:
$\fn{Utilidade}(s, O) = -\fn{Utilidade}(s, X)$.\newline
\textbf{Informação perfeita:} sim, o tabuleiro inteiro é visível para os dois jogadores.\newline
\textbf{Determinístico:} sim, $\fn{Resultado}(s, a)$ é único e não há sorteio.\newline
\textbf{Finito:} sim, cada jogada ocupa uma das 42 casas, então toda partida acaba em no máximo 42 jogadas. \\
\bottomrule
\end{tabularx}

% ---------------------------------------------------------------------
\secao{2. Formalização (ID2.1)}
\begin{tabularx}{\linewidth}{@{}LY@{}}
\toprule
\campo{Estado e $s_0$}{incluindo de quem é a vez. Sem isso, \fn{Jogador}(s) não é função do estado} &
$s = (x, o, \mathit{alturas}, \mathit{vez})$. $x$ e $o$ são bitboards de 49 bits (7 colunas × 7 bits, com 1 bit
sentinela por coluna) com as casas ocupadas por X e por O; $\mathit{alturas}[c] \in \{0..6\}$ é o número de peças na
coluna $c$; $\mathit{vez} \in \{X, O\}$ é quem joga. $\fn{Jogador}(s) = s.\mathit{vez}$.\newline
$s_0$ = tabuleiro vazio, $\mathit{alturas} = (0, \dots, 0)$, $\mathit{vez} = X$ (MAX começa). \\
\midrule
\campo{\fn{Ações} e \fn{Resultado}}{jogadas legais e o modelo de transição} &
$\fn{Ações}(s) = \{ c \in \{0..6\} : \mathit{alturas}[c] < 6 \}$, ou $\emptyset$ se $s$ é terminal. Logo
$|\fn{Ações}(s)| \le 7$.\newline
$\fn{Resultado}(s, c)$: coloca a peça de $\fn{Jogador}(s)$ na casa (coluna $c$, linha $\mathit{alturas}[c]$), incrementa
$\mathit{alturas}[c]$ e passa a vez. Jogar em coluna cheia é erro. \\
\midrule
\campo{\fn{Teste-Terminal}}{quando o jogo acaba} &
Verdadeiro se X ou O tem 4 em linha ou se as 42 casas estão ocupadas. O alinhamento é testado no bitboard: para cada
uma das 4 direções, $b \mathbin{\&} (b{\gg}d) \mathbin{\&} (b{\gg}2d) \mathbin{\&} (b{\gg}3d) \ne 0$. \\
\midrule
\campo{\fn{Utilidade}(s, MAX)}{sempre do ponto de vista de MAX. Um único número, não dois} &
$\fn{Utilidade}(s, \text{MAX}) = +(1000 - n)$ se X venceu, $-(1000 - n)$ se O venceu e 0 no empate, onde $n$ é o número
de peças no tabuleiro ($n \le 42$).\newline
É um único número; $\fn{Utilidade}(s, \text{MIN}) = -\fn{Utilidade}(s, \text{MAX})$. O desconto por $n$ depende só do
estado e faz a vitória mais rápida valer mais, o que impede o agente de adiar uma vitória que já viu. Toda vitória
vale pelo menos 958 em módulo e a avaliação heurística é limitada a $\pm 900$, então um corte nunca se confunde com
um terminal. \\
\midrule
\campo{$b$ e $m$ estimados}{fator de ramificação médio e profundidade típica, com a conta de $b^m$} &
$b$: 7 no início, caindo quando colunas enchem. Média medida nos estados em que o agente jogou: $b \approx \dado{b-medio-1}$.\newline
$m$: no máximo 42 jogadas. Duração média medida: \dado{m-medio} jogadas (\dado{d6-durmedia} contra AlfaBeta-d6).\newline
Conta: o limite superior é $b^m = 7^{42} \approx \dado{bm-max}$ folhas; com os valores medidos,
$\dado{b-medio-1}^{\,\dado{m-medio}} \approx \dado{bm-medido}$. As posições legais distintas são $\approx 4{,}5 \times 10^{12}$
(contagem de Tromp). Mesmo a $10^6$ nós/s a árvore completa é inviável, por isso usamos alfa-beta com corte de
profundidade. \\
\bottomrule
\end{tabularx}

% ---------------------------------------------------------------------
\secao{3. Implementação (ID2.2)}
\begin{tabularx}{\linewidth}{@{}LY@{}}
\toprule
\campo{Ordenação de jogadas}{que heurística barata você usa para tentar examinar primeiro a melhor jogada} &
Centro primeiro: as colunas são examinadas na ordem 3, 2, 4, 1, 5, 0, 6 (tabela fixa, custo zero). A coluna central
participa de \dado{inclusoes-centro} das \dado{inclusoes} inclusões casa-janela (a de borda, de \dado{inclusoes-borda}),
então as jogadas centrais costumam ser as melhores e provocam podas cedo.\newline
Medido em \dado{ord-posicoes} posições com $d = 8$: média de \dado{ord-com-media} nós com ordenação contra
\dado{ord-sem-media} sem (medianas \dado{ord-com-mediana} e \dado{ord-sem-mediana}), com o mesmo valor na raiz. Em
\dado{ord-pior} das \dado{ord-posicoes} posições a ordenação custou mais nós. \\
\midrule
\campo{Função de avaliação}{as características e os pesos, com a justificativa de cada um} &
$\fn{Aval}(s) = \sum_{\text{janelas}} [5\cdot(\text{3 de X + 1 vazia}) + 2\cdot(\text{2 de X + 2 vazias})] - (\text{o mesmo para O})
+ 3\cdot(\text{peças de X} - \text{peças de O na coluna central})$, limitada a $\pm 900$.\newline
As \dado{janelas} janelas são todos os segmentos de 4 casas em linha. Uma janela com peças dos dois lados vale 0, porque
ninguém pode mais vencer por ela.\newline
Peso 5, 3 abertas: ameaça a uma jogada da vitória. Peso 2, 2 abertas: ameaça em construção, ainda faltam 2 jogadas.
Peso 3 por peça no centro: a casa central está em \dado{janelas-max-casa} janelas, contra \dado{janelas-canto} de um
canto, então o centro dá controle posicional. Os pesos foram escolhidos à mão, sem ajuste fino, e são os mesmos para
todos os agentes de busca. \\
\midrule
\campo{Mecanismo de parada}{profundidade fixa, aprofundamento iterativo com relógio, quiescência} &
Profundidade fixa $d$ (agente avaliado $d = 8$; adversários $d = 2$, 4 e 6), cortando também nos terminais.\newline
Preferimos profundidade fixa a relógio porque o custo é medido em nós: com relógio, a quantidade de nós e as jogadas
dependeriam da máquina e o experimento não seria reprodutível. Não há busca de quiescência, então o efeito horizonte é
possível (limitação declarada). \\
\bottomrule
\end{tabularx}

% ---------------------------------------------------------------------
\secao{4. Protocolo experimental (ID2.3)}
\begin{tabularx}{\linewidth}{@{}LY@{}}
\toprule
\campo{Adversários}{mínimo obrigatório: um agente aleatório e o seu próprio agente com profundidade menor} &
(1) Aleatório: escolha uniforme em $\fn{Ações}(s)$, com semente fixa. É a linha de base trivial.\newline
(2) O próprio agente com profundidade menor: AlfaBeta-d2, d4 e d6, com o mesmo código, a mesma \fn{Aval}, os mesmos pesos
e a mesma ordenação. Só a profundidade muda, então a calibração é simétrica. \\
\midrule
\campo{Unidade de custo}{nós visitados. Não gerações, não segundos} &
Nós visitados: cada chamada à função recursiva de busca conta 1, incluindo a raiz, as folhas cortadas por profundidade
e os terminais.\newline
Para cada lance do agente avaliado reportamos média, desvio, mediana, p95 e máximo. O tempo de parede só é registrado
como informação e não entra nas comparações. \\
\midrule
\campo{Partidas e semente}{quantas partidas, quem começa, como as aberturas são variadas} &
vs AlfaBeta-d2/d4/d6: as 2 primeiras jogadas são forçadas e percorrem as $7 \times 7 = 49$ aberturas, cada uma jogada
2 vezes trocando quem começa: 98 partidas por confronto. Não há sorteio; empates entre jogadas são desfeitos pela ordem
de exame.\newline
vs Aleatório: 100 partidas, com semente do adversário igual ao número da partida (0 a 99); o agente avaliado começa
nas partidas pares.\newline
Total de \dado{total-partidas} partidas. A amostra do teste de ordenação usa semente 42. Tudo é reproduzível com
\texttt{python3 experimento.py}. \\
\bottomrule
\end{tabularx}

\aviso{Uma armadilha específica de jogos:}{dois agentes determinísticos jogam sempre a MESMA partida. Cem partidas
seriam uma partida repetida cem vezes. Varie a abertura por sorteio, ou alterne quem começa, e declare isso no protocolo.}

% ---------------------------------------------------------------------
\secao{5. Verificação antes de sair da aula}
{\small Marque cada item e mostre ao professor. O entregável do bloco não é a ficha preenchida, é o código rodando.}

\begin{tabularx}{\linewidth}{@{}cY@{}}
\toprule
\boldmath$\checkmark$ & \textbf{Item} \\
\midrule
\ok & O jogo passa nas cinco verificações da seção 1 \\
\ok & O estado inclui de quem é a vez \\
\ok & A utilidade é um único número, do ponto de vista de MAX \\
\ok & Minimax roda em profundidade completa em uma instância pequena do seu jogo
      {\small (\texttt{verificacao.py}: Lig 4 4×4 com $K = 3$, a partir de $s_0$: valor \dado{t1-mm-valor}, \dado{t1-mm-nos} nós)} \\
\ok & Alfa-beta devolve exatamente o mesmo valor que Minimax nessa instância
      {\small (4×4 $K=3$: \dado{t1-ab-valor} com \dado{t1-ab-nos} nós; também 200 posições 3×4, 100 finais 6×7 em profundidade
      completa e 100 posições 6×7 com $d = 4$: zero divergências)} \\
\bottomrule
\end{tabularx}

\aviso{Sobre o último item, e ele é o mais importante:}{a igualdade de valor entre Minimax e alfa-beta em profundidade
completa é a única evidência de corretude que o Teorema 1 da Aula 7 autoriza exigir, e é barata de produzir. Sem ela
não há como distinguir um alfa-beta correto de um que poda demais. Este item vale nota.}

% ---------------------------------------------------------------------
\secao{6. Critérios de avaliação}
\begin{tabularx}{\linewidth}{@{}Ycc@{}}
\toprule
\textbf{Critério} & \textbf{Peso} & \textbf{Indicador} \\
\midrule
Formalização do jogo: os seis componentes especificados e justificados & 1,0 & ID2.1 \\
Alfa-beta correto, verificado por igualdade de valor com Minimax & 1,0 & ID2.2 \\
Protocolo reprodutível, com unidade de custo e calibração simétricas & 1,0 & ID2.3 \\
Análise dos resultados, com distribuição reportada e garantias ausentes declaradas & 1,0 & ID2.3 \\
\bottomrule
\end{tabularx}

% ---------------------------------------------------------------------
\secao{7. As três formas de perder nota por comparação injusta}
{\small Elas valem nota, e estão anunciadas de antemão.}

\begin{tabularx}{\linewidth}{@{}>{\raggedright\arraybackslash}p{0.2\linewidth}Y>{\raggedright\arraybackslash}p{0.28\linewidth}@{}}
\toprule
\textbf{A falha} & \textbf{Como se manifesta} & \textbf{Como evitar} \\
\midrule
Unidade de custo incomparável & Comparar gerações com iterações, ou qualquer coisa com segundos. Uma geração com
população 50 custa 50 avaliações; uma iteração de Simulated Annealing custa 1; um lance de busca custa milhares de nós
& Medir tudo na mesma unidade, declarada no relatório \\
Calibração assimétrica & Varrer dezenas de configurações de um método e comparar contra o outro nos valores padrão.
Isso mede a diferença entre calibrar e não calibrar & Mesmo esforço de calibração dos dois lados, declarado no relatório \\
Melhor caso, sem linha de base & Reportar a melhor de N execuções. Isso mede a sorte, não o método & Reportar média,
pior caso e taxa de sucesso. Sempre incluir a linha de base trivial \\
\bottomrule
\end{tabularx}

% ---------------------------------------------------------------------
\secao{8. Armadilhas na escolha do jogo}
\begin{tabularx}{\linewidth}{@{}>{\raggedright\arraybackslash}p{0.18\linewidth}Y>{\raggedright\arraybackslash}p{0.32\linewidth}@{}}
\toprule
\textbf{Escolha} & \textbf{O que dá errado} & \textbf{Redirecionamento} \\
\midrule
Jogo da velha & Minimax puro resolve inteiro, então alfa-beta não é necessário e o corte de profundidade não faz sentido
& Aceito apenas como instância de verificação, nunca como o jogo do trabalho \\
Xadrez completo & Implementar as regras consome o prazo inteiro e sobra zero tempo para o experimento & Escolher um
subconjunto, como finais de rei e torre, ou trocar de jogo \\
Jogo com sorteio de dados & Deixa de ser determinístico, e o valor de um nó vira uma esperança & Fora do escopo do
RA2. Trocar de jogo \\
Jogo com informação oculta & Estratégias puras deixam de bastar, e a solução exige aleatorização & Fora do escopo do
RA2. Trocar de jogo \\
\bottomrule
\end{tabularx}

% ---------------------------------------------------------------------
\secao{Roteiro da apresentação}
{\small Em 27/10, cada equipe tem 12 minutos, com 3 minutos de arguição para apresentar o TDE2 ao professor. A
apresentação é avaliada pelo mesmo instrumento do relatório, e não separadamente.}

\begin{tabularx}{\linewidth}{@{}cY@{}}
\toprule
\textbf{Tempo} & \textbf{O que mostrar} \\
\midrule
1 min & O problema e por que ele é adequado ao método. Sem introdução histórica \\
3 min & A modelagem: representação, o tamanho do espaço com a conta, e o que a representação torna impossível por construção \\
3 min & Uma execução ao vivo ou gravada, com o contador de custo visível na tela \\
3 min & A tabela de comparação contra as linhas de base, na mesma unidade e com a distribuição, não só o melhor caso \\
2 min & O que não funcionou, e as garantias que o seu sistema não tem \\
\bottomrule
\end{tabularx}

\aviso{O item que mais pesa e que as equipes costumam pular:}{o último. Uma apresentação que só mostra o que deu certo
perde pontos no indicador de avaliação de resultados. Declarar explicitamente uma limitação medida vale mais que um
resultado bonito sem ressalva.}

\end{document}
